Me preparo
L1.   El espacio geográfico: interrelaciones sociedad-naturaleza
L2.   Categorías de análisis del espacio geográfico
L3.   El estudio del espacio geográfico
L4.   Representaciones cartográficas
L5.   Interpretación y representación de la Información geográfica
L6.   Tecnologías de la información geográfica
L7.   Las placas tectónicas y su relación con la sismicidad y el vulcanismo
L8.   Las placas tectónicas y su relación con el relieve
L9.   El suelo y su importancia
L10.  Aguas continentales en México y el mundo
L11.  Las cuencas hidrográficas y el desarrollo económico en México
L12.  Aguas oceánicas y la importancia de los litorales de México


F1.  Los pueblos indígenas y su relación con la naturaleza
F2.  ¿Qué muestran los mapas?
F3.  Un paso adelante, menos
F4.  Planta sobre planta, caminemos
F5.  Cuencas y disponibilidad de agua limpia
F6.  Océanos en peligro ¡Paremos lacontaminación!

     